\chapter{Financial and Mathematical Background}

\section{European Call Options}

A European call option gives its holder the right, but not the obligation, to buy an underlying asset at a fixed strike price $K$ at maturity $T$.
If the terminal asset price is $S_T$, the payoff is:

\begin{equation}
    \max(S_T - K, 0).
\end{equation}

The option value depends on the current asset price, the strike price, time to maturity, the risk-free rate, and the volatility of the underlying asset.
General arbitrage-pricing principles for derivative securities can be introduced in discrete-time models before moving to continuous-time diffusion models \cite{shreve2004stochastic1,pascucci2009finanza}.
In this project, only European call options are considered.

\section{Stochastic Differential Equations}

Continuous-time asset price models are often described by stochastic differential equations.
An SDE combines a deterministic drift component with a random diffusion component driven by Brownian motion.
This framework is standard in financial modeling because asset prices evolve continuously but are subject to uncertainty \cite{shreve2004stochastic}.

\section{Geometric Brownian Motion}

In the Black-Scholes model, the underlying asset follows a geometric Brownian motion under the risk-neutral measure:

\begin{equation}
    dS_t = r S_t dt + \sigma S_t dW_t,
\end{equation}

where $S_t$ is the asset price at time $t$, $r$ is the risk-free interest rate, $\sigma$ is the constant volatility, and $W_t$ is a standard Brownian motion.

The solution at maturity $T$ is:

\begin{equation}
    S_T = S_0 \exp\left[\left(r - \frac{1}{2}\sigma^2\right)T + \sigma \sqrt{T} Z\right],
    \qquad Z \sim \mathcal{N}(0,1).
\end{equation}

\section{Black-Scholes Model}

The Black-Scholes model assumes frictionless markets, constant volatility, a constant risk-free rate, no arbitrage, and continuous trading.
Under these assumptions, the price of a European option can be derived from risk-neutral valuation.

For a European call, the value is the discounted expected payoff:

\begin{equation}
    C = e^{-rT}\mathbb{E}^{\mathbb{Q}}\left[\max(S_T - K, 0)\right].
\end{equation}

\section{Black-Scholes Formula}

The analytical Black-Scholes price of a European call option is:

\begin{equation}
    C_{BS} = S_0 N(d_1) - K e^{-rT} N(d_2),
\end{equation}

where $N(\cdot)$ is the cumulative distribution function of a standard normal random variable and:

\begin{align}
    d_1 & = \frac{\ln(S_0/K) + \left(r + \frac{1}{2}\sigma^2\right)T}{\sigma\sqrt{T}}, \\
    d_2 & = d_1 - \sigma\sqrt{T}.
\end{align}

In this project, $C_{BS}$ is used both as the supervised learning target and as the reference value for evaluation.
This is what makes the experiment controlled: the true target function is known analytically.

\section{Monte Carlo Pricing}

Monte Carlo simulation approximates the risk-neutral expectation by generating many realizations of $S_T$:

\begin{equation}
    \hat{C}_{MC} =
    e^{-rT}
    \frac{1}{M}
    \sum_{i=1}^{M}
    \max(S_T^{(i)} - K, 0).
\end{equation}

For sufficiently large $M$, $\hat{C}_{MC}$ converges to $C_{BS}$.
In this project, Monte Carlo is used as a numerical benchmark and as a comparison point for the neural network.
Since the terminal distribution of $S_T$ is known exactly under Black-Scholes, the implementation simulates terminal prices directly rather than discretizing the full path.
